\section{Data, design and notation}\label{sec:data}

\subsection{Three datasets}

The problem supplies three kinds of evidence, summarised in Table~\ref{tab:data}.

\begin{table}[htbp]
  \centering
  \caption{Datasets used in this paper.}\label{tab:data}
  \small
  \begin{tabular}{llll}
    \toprule
    Source & Content & Shape & Used in \\
    \midrule
    Data01 & acoustic measurements of sea-surface calls & 108 calls $\times$ 11 fields & Problem 1 \\
    Data02 & playback trials, paired call/control & 278 rows = 139 pairs $\times$ 40 fields & Problem 3, 4 \\
    Data03--11 & propagation correlations at 7 distances & 9 files $\times$ 300 rows $\times$ 7 distances & Problem 2, 4 \\
    \bottomrule
  \end{tabular}
\end{table}

The acoustic file contains 44 flat, 40 modulated and 24 two-voice calls. Two fields are
\emph{structurally} missing rather than lost: the second fundamental $F_{02}$ exists only for
two-voice calls (24/24 present, 0/84 elsewhere), and the pulse rate is present for modulated
(34/40) and two-voice (20/24) calls but never for flat calls (0/44). This structure is a
property of the call types, and treating it as an ordinary covariate would leak the label into
the classifier; Section~\ref{sec:q1} quantifies that leakage.

\subsection{The paired playback design}\label{sec:paired}

Each of the 139 experiments (\texttt{expID}) contains exactly two rows: one natural call and one
control, played on the same date and in the same behavioural context. This is a matched-pair
design, and it is the single most important structural fact in the dataset. Three consequences
follow.

First, the comparison of interest is within-pair. A test that pools the 139 call rows and the
139 control rows discards the matching and attenuates the effect: the pooled odds ratio is 12.4,
whereas the paired McNemar odds ratio is 37.0 (Section~\ref{sec:q3-overall}).

Second, matching is complete at the level of date and context but not at the level of distance.
All 139 pairs share a date and a context; only 49 pairs share a distance, and the remaining 90
differ by between $-150$ and $+120$ m. Pairing is therefore a group-level device, not a
distance-matched one. We use the pair as the clustering unit (generalised estimating equations
with an exchangeable working correlation) and report an equal-distance subset as a sensitivity
analysis rather than as the primary estimate.

Third, the 139 call trials reuse 28 distinct playback tracks (plus one pseudo-track for the
control), with a mean of 4.96 trials per track. Trials are therefore not independent at the level
of the stimulus either. The track-level intraclass correlation of response is 0.079, which is
small but non-zero; we treat it as a sensitivity dimension rather than as a second random effect
in the primary model, because the number of tracks per call type is small (11, 7 and 10).

\subsection{Derived outcomes and definitional sensitivity}\label{sec:derived}

The response indicators in Data02 are derived quantities, not independent observations. The
identities hold for all 278 rows without exception:
\begin{align}
  \text{HeadsUpYesNo} &= \mathbf{1}\{\text{HeadsUpDuration}>0\}, \\
  \text{ApproachYesNo} &= \mathbf{1}\{\text{ApproachDuration}>0\}, \\
  \text{VocalYesNo} &= \mathbf{1}\{\text{VocalNcalls}>\text{VocalNcallsBefore}\}, \\
  \text{ResponseYesNo} &= \text{HeadsUp} \lor \text{Vocal} \lor \text{Approach}.
\end{align}

The third identity deserves emphasis. \texttt{VocalYesNo} measures a \emph{net increase} over the
pre-playback baseline, not the occurrence of any vocalisation. Four rows therefore have
\texttt{VocalYesNo}$=0$ while \texttt{VocalNcalls}$>0$: the bird vocalised after playback, but
less than it had before. These rows are not data errors, and we do not delete them. To show that
the conclusions do not depend on this convention, we recompute the response counts under three
alternative definitions (Table~\ref{tab:defsens}).

\begin{table}[htbp]
  \centering
  \caption{Definitional sensitivity of the response counts (all 278 rows, calls and controls
  combined). S0 is the primary definition.}\label{tab:defsens}
  \small
  \begin{tabular}{llrrr}
    \toprule
    Scenario & Definition & Vocal events & Any response (278 rows) & of which control \\
    \midrule
    S0 & net increase (primary) & 81 & 108 & 18 \\
    S1 & any post-playback vocalisation & 85 & 111 & 20 \\
    S2 & approach attempt (latency recorded) & 81 & 109 & 18 \\
    S3 & S1 and S2 combined & 85 & 112 & 20 \\
    \bottomrule
  \end{tabular}
\end{table}

Across the three alternatives the total number of any-response trials (calls and controls
combined) moves between 108 and 112, of which 18--20 are controls, so the qualitative conclusions
of Section~\ref{sec:q3} are not an artefact of the definition.

\subsection{The nested structure of the propagation data}

Each of the nine propagation files contains 300 rows and seven columns, one per distance. The
rows are not 300 independent trials. They form 30 contiguous blocks of ten rows, and each block
corresponds to one call-level unit measured at all seven distances. A block indicator alone
explains 55.7\% (two-voice), 72.2\% (flat) and 80.2\% (modulated) of the row-level variance in
fidelity; the median design effect is 6.46, which corresponds to an effective sample size of
about 46 per distance rather than 300. Any interval computed by resampling rows independently is
too narrow by a factor of about 2.5. All uncertainty statements about the CEI radii in
Section~\ref{sec:q4} therefore use a \emph{block} bootstrap that resamples the 30 blocks, not
the rows \cite{ref7}. The Q2 threshold radii in Table~\ref{tab:q2-ranges} use a record-level bootstrap and
should be read as lower bounds; the model-family range is the honest summary.

Within a call type, the missing-value pattern is identical across the three metrics, which
confirms that the three columns describe the same trials and can be averaged row-wise. Across
call types the rows cannot be matched, because no trial identifier is shared; we never assume that
row $i$ of the flat file and row $i$ of the modulated file are the same event.

\subsection{Assumptions}

Four assumptions carry the analysis, and each is stated with the evidence that supports it.

\begin{description}[leftmargin=1.2em,itemsep=2pt]
  \item[A1] \textbf{The three correlations measure one latent fidelity.} Supported by the strong
  correlation between the spectrogram and amplitude-spectrum metrics ($r=0.81$--$0.97$ per call
  type). We nevertheless treat the equal weighting as a choice, not a fact, and vary it in
  Section~\ref{sec:q4-sens}.
  \item[A2] \textbf{The playback pair is exchangeable within an experiment.} Justified by design
  (same date, context and, in 49 pairs, same distance); relaxed by the equal-distance subset.
  \item[A3] \textbf{Censoring is non-informative at $W=60$ s.} The observation window is not
  recorded in the data; 60 s is the upper bound of the duration field and exceeds the largest
  observed latency (47.8 s). We vary $W$ over 48, 90 and 120 s.
  \item[A4] \textbf{Fidelity and recruitment are multiplicatively separable.} This is a modelling
  choice with an ecological justification (two necessary conditions), not a testable law; we say
  so explicitly in Section~\ref{sec:q4}.
\end{description}
